\documentclass[11pt,a4paper]{article}
\usepackage[margin=24mm]{geometry}
\usepackage{amsmath,amssymb,bm,booktabs,longtable,array,hyperref,xcolor}
\hypersetup{colorlinks=true,linkcolor=teal,urlcolor=teal,citecolor=teal,pdftitle={Einstein--Cartan spherical collapse: derivation and readiness audit}}
\newcommand{\dd}{\mathrm d}
\newcommand{\Dt}{D_t}
\newcommand{\Dr}{D_r}
\newcommand{\eps}{\epsilon}
\newcommand{\kap}{\kappa}
\newcommand{\Tr}{\mathcal T}
\newcommand{\Rsc}{\mathcal R}
\newcommand{\Lie}{\mathcal L}
\newcommand{\status}[1]{\textbf{#1}}
\setlength{\emergencystretch}{3em}
\title{Einstein--Cartan spherical collapse\\Mathematical derivation, compatibility and validation audit}
\author{EC Collapse Laboratory research record\\\small Repository: \texttt{pdamkier-del/skovlab}}
\date{10 November 2026}
\begin{document}
\maketitle
\begin{abstract}
This is a derivation and implementation-readiness audit, not a spherical-collapse simulation. We derive the spherical metric equations, their constraints, mass laws, conservation laws and FLRW reductions. We distinguish a spinless fluid, a conditionally averaged Hashemi-normalized spin fluid and a published unaveraged Weyssenhoff formulation. The central obstruction is that the natural local extension of conserved-spin ordinary dust has negative acoustic speed squared on its low-spin branch. Exact FLRW symmetry also excludes a nonzero coherent Frenkel spin tensor. These facts prohibit an unqualified promotion of the existing homogeneous bounce model to a general spherical PDE. No application, existing numerical solver or existing model equations have been changed. No new exterior matching construction is claimed.
\end{abstract}
\tableofcontents
\section{Scope, conventions and variational equations}
\subsection{What is verified and what is conditional}
The spherical geometry and GR perfect-fluid equations below are derived and independently checked. The averaged EC source is a \emph{specified effective matter model}, not a consequence of symmetry alone. Its local extension is algebraically closed under explicit transport assumptions, but not necessarily a well-posed Cauchy problem. For full Weyssenhoff matter we select the established Obukhov--Korotky/Br\'echet formulation \cite{OK,BHL}; we do not manufacture a new microscopic matter action. Its coherent spherically symmetric branch is distinguished from an unclosed hierarchy of averaged spin correlations. ``Verified'' below means algebraically checked identities and analytical limits, not a general existence theorem, a physical validation of a matter closure or proof of a complete black-hole spacetime.

\subsection{Connection, curvature and units}
Take $g_{ab}=\operatorname{diag}(-,+,+,+)$ in an orthonormal frame, $c=1$, $\kap=8\pi G$, and no cosmological constant or Holst term. The connection convention is
\begin{align}
 \widetilde\nabla_\mu V^\rho&=\partial_\mu V^\rho+\widetilde\Gamma^\rho{}_{\mu\nu}V^\nu, &
 \Tr^\rho{}_{\mu\nu}&=2\widetilde\Gamma^\rho{}_{[\mu\nu]},\\
 K^\rho{}_{\mu\nu}&=\widetilde\Gamma^\rho{}_{\mu\nu}-\{^\rho{}_{\mu\nu}\}, &
 K_{abc}&=\tfrac12(\Tr_{abc}-\Tr_{bca}+\Tr_{cab}).\label{eq:contortion}
\end{align}
The derivative index is the \emph{first lower} connection index. Metric compatibility gives $K_{abc}=-K_{cba}$. Antisymmetrization and symmetrization have weight $1/2$. Our curvature convention is
\begin{equation}
 \widetilde R^\rho{}_{\sigma\mu\nu}=\partial_\mu\widetilde\Gamma^\rho{}_{\nu\sigma}-\partial_\nu\widetilde\Gamma^\rho{}_{\mu\sigma}
 +\widetilde\Gamma^\rho{}_{\mu\lambda}\widetilde\Gamma^\lambda{}_{\nu\sigma}
 -\widetilde\Gamma^\rho{}_{\nu\lambda}\widetilde\Gamma^\lambda{}_{\mu\sigma},
 \qquad \widetilde R_{\sigma\nu}=\widetilde R^\rho{}_{\sigma\rho\nu}.
\end{equation}
Untilded curvature and $\nabla$ always use the Levi--Civita connection. The metric Ricci scalar is $\Rsc$; it is not assumed equal to $\widetilde R$. Coordinates $t,r$ can be measured in length units, $R$ is an areal length and $\Phi,\Lambda$ are dimensionless. In $R=a(t)r$ with dimensional $r$, $a$ is dimensionless and $k$ has units of inverse length squared. The simulator instead normalizes all lengths to $L_*$ and uses a dimensionless comoving radius and $k=0,\pm1$; a coordinate rescaling relates the two conventions. The equations are tensorial under that relabelling.

\subsection{Gravitational variation and torsion elimination}
Start from
\begin{equation}\label{eq:action}
 I=\frac1{2\kap}\int\sqrt{-g}\,\widetilde R\,\dd^4x+I_m[e,\widetilde\omega,\psi].
\end{equation}
An orthonormal coframe $e^a$, metric-compatible $\widetilde\omega^{ab}=-\widetilde\omega^{ba}$ and matter fields are independent. Fix orientation $\eta=e^0\wedge e^1\wedge e^2\wedge e^3$, $\epsilon_{0123}=+1$, and define the \emph{operational} spin source by
\begin{equation}
 \delta I_m=\int\delta e^a\wedge\Sigma_a-\frac12\int\delta\widetilde\omega^{ab}\wedge\tau_{ab}+\int E_\psi\delta\psi,
 \qquad \tau_{ab}=\tau^\mu{}_{ab}\,\iota_{\partial_\mu}\eta.
\end{equation}
In these conventions the gravitational integrand is $(4\kap)^{-1}\epsilon_{abcd}e^a\wedge e^b\wedge\widetilde{\mathcal R}^{cd}$. Use $\delta\mathcal R^{ab}=\widetilde D\delta\widetilde\omega^{ab}$ and integrate once by parts. The independent equations are
\begin{align}
 \epsilon_{abcd}e^b\wedge\widetilde{\mathcal R}^{cd}&=-2\kap\Sigma_a,\\
 \epsilon_{abcd}\Tr^a\wedge e^b&=\kap\tau_{cd}, & E_\psi&=0.\label{eq:cartanform}
\end{align}
Variations have compact support in this bulk calculation. A complete boundary variational principle requires appropriate boundary/interface terms, not supplied by dropping a total derivative. The tetrad equation is the torsionful Einstein equation in this source convention. The second equation is algebraic in torsion. It is not a wave equation. Local Lorentz variation also gives the matter spin identity $\widetilde D\tau_{ab}+2e_{[a}\wedge\Sigma_{b]}=0$ on the matter equations. The bare canonical stress need not be symmetric or Levi--Civita conserved.

For an ideal convective source $\tau^\lambda{}_{\mu\nu}=u^\lambda s_{\mu\nu}$ with $s_{\mu\nu}u^\nu=0$, Eq.~\eqref{eq:cartanform} has zero torsion trace and yields $\Tr^\lambda{}_{\mu\nu}=\kap u^\lambda s_{\mu\nu}$. This fixes the normalization of $s$ \emph{here}; changing the sign in the definition of the spin variation changes its linear-sign dictionary. Let $s^2=\frac12s_{\mu\nu}s^{\mu\nu}\geq0$.

There is a useful independent coefficient check. For a minimal matter coupling linear in contortion, in a local rest frame the quadratic gravitational contortion term and linear matter term evaluated at the Cartan solution are
\begin{equation}
 \Delta\widetilde R=-\frac{\kap^2}{2}s^2,\qquad
 \mathcal L_{g,K}=-\frac\kap4s^2,\qquad
 \mathcal L_{m,K}=-\tfrac12u^\mu s_{ab}K^{ab}{}_{\mu}=\frac\kap2s^2.
\end{equation}
The sum is $+\kap s^2/4$. The divergence in the scalar decomposition vanishes for the trace-free convective source. This can also be understood by eliminating a quadratic auxiliary variable: at its stationary value the linear term is minus twice the quadratic term. The check uses three arbitrary spatial spin components, verifies Eq.~\eqref{eq:contortion} and metric compatibility, and is reproduced symbolically. It does \emph{not} replace a constitutive matter action: spin transport, internal energy and averaging must still be specified. Nonminimal torsion couplings would invalidate this coefficient argument.

\subsection{A source dictionary, not a blind coefficient transfer}
Hashemi et al. define half-antisymmetric connection torsion and use a matter-source variation with an explicit factor two; their averaged source subtracts $\kap S_H^2/4$ \cite{Hashemi}. Br\'echet et al. use $Q=\Gamma_{[\mu\nu]}$, derivative index last, $\widetilde\Gamma=\Gamma-K_B$, and subtract $\kap b^2$ in their spin variable $b_{\mu\nu}$ \cite{BHL}. With the same geometric connection, our torsion is $\Tr^\rho{}_{\mu\nu}=-2Q_B^\rho{}_{\mu\nu}$, and $K^\rho{}_{\mu\nu}=-(K_B)^\rho{}_{\nu\mu}$. Their $Q_B=\kap u b$ corresponds to $s=-2b$ in the operational source convention above. Thus the \emph{scalar} effective correction can be aligned as
\begin{equation}\label{eq:dictionary}
 q\equiv\kap S_H^2/4=\kap s^2/4=\kap b^2,\qquad S_H^2=s^2=4b^2.
\end{equation}
This alignment requires the same ideal source and averaging assumptions. It is not a licence to identify the same numerical tensor called ``spin'' in different actions. Linear-spin terms retain their own signs and source definitions. In particular, a scalar matching does not prove equivalence of microscopic models or averaging prescriptions. The existing FLRW input remains $\overline{S_H^2}=\kap^2L_*^2 S_H^2$, and $\bar q=\overline{S_H^2}/4$.

\section{Spherical geometry: derivation before matter substitution}
Write $A=e^\Phi$, $B=e^\Lambda$, and
\begin{equation}
 g=h_{ab}\dd x^a\dd x^b+R^2\gamma_{ij}\dd y^i\dd y^j,
 \quad h_{ab}=\operatorname{diag}(-A^2,B^2),\quad \gamma=\dd\theta^2+\sin^2\theta\dd\phi^2.
\end{equation}
Dots and primes in component expressions mean coordinate $\partial_t,\partial_r$. Set $X=h^{ab}R_{,a}R_{,b}$ and $H_{ab}=\nabla^{(h)}_a\nabla^{(h)}_bR$. The mixed Christoffel blocks, besides those intrinsic to $h$ and the unit sphere, are
\begin{equation}
 \Gamma^a{}_{ij}=-R R^{,a}\gamma_{ij},\qquad
 \Gamma^i{}_{aj}=\frac{R_{,a}}R\delta^i_j.
\end{equation}
Inserting these blocks in the Ricci definition gives
\begin{align}
 R^{(4)}_{ab}&=R^{(h)}_{ab}-2H_{ab}/R,\\
 R^{(4)}_{ij}&=[1-X-R\Box_hR]\gamma_{ij},\\
 \Rsc&=R_h+2(1-X)/R^2-4\Box_hR/R.
\end{align}
Because $R^{(h)}_{ab}=\tfrac12R_hh_{ab}$ in two dimensions, the independent Einstein blocks are
\begin{equation}\label{eq:warped}
 G_{ab}=-2H_{ab}/R+\left[2\Box_hR/R+(X-1)/R^2\right]h_{ab},\qquad
 G^\theta{}_\theta=\Box_hR/R-R_h/2.
\end{equation}
This is the documented geometric derivation. The CAS check instead constructs all four-dimensional Christoffels and Ricci components directly from the metric, without using Eq.~\eqref{eq:warped}.

\subsection{Explicit independent components}
Define $E_0=G_{tt}/A^2$, $E_1=G_{rr}/B^2$, $E_2=G_{\theta\theta}/R^2$. Direct evaluation of the Hessians gives
\begin{align}
 E_0={}&e^{-2\Phi}\left(\frac{2\dot R\dot\Lambda}{R}+\frac{\dot R^2}{R^2}\right)
 -e^{-2\Lambda}\left(\frac{2R''}{R}+\frac{R'^2}{R^2}-\frac{2\Lambda'R'}R\right)+\frac1{R^2},\label{eq:E0}\\
 E_1={}&-e^{-2\Phi}\left(\frac{2\ddot R}R+\frac{\dot R^2}{R^2}-\frac{2\dot\Phi\dot R}R\right)
 +e^{-2\Lambda}\left(\frac{R'^2}{R^2}+\frac{2\Phi'R'}R\right)-\frac1{R^2},\label{eq:E1}\\
 G_{tr}={}&-\frac2R(\dot R'-\Phi'\dot R-\dot\Lambda R'),\label{eq:Etr}\\
 E_2={}&-e^{-2\Phi}\left[\ddot\Lambda+\dot\Lambda^2-\dot\Phi\dot\Lambda+
 \frac{\ddot R+(\dot\Lambda-\dot\Phi)\dot R}{R}\right]\nonumber\\
 &+e^{-2\Lambda}\left[\Phi''+\Phi'^2-\Phi'\Lambda'+
 \frac{R''+(\Phi'-\Lambda')R'}R\right].\label{eq:E2}
\end{align}
$G_{\phi\phi}=\sin^2\theta G_{\theta\theta}$; all other off-diagonal components vanish. These are independent component equations, not four independent propagating gravitational degrees of freedom: Bianchi identities relate their constraint propagation and evolution.

\subsection{General spherical source and conservation}
Choose a comoving, vorticity-free fluid congruence $u=A^{-1}\partial_t$ and outward unit radial vector $n=B^{-1}\partial_r$. A general symmetric spherical source can be decomposed as
\begin{equation}\label{eq:stress}
 T_{\mu\nu}=\eps u_\mu u_\nu+P_r n_\mu n_\nu+P_\perp(g_{\mu\nu}+u_\mu u_\nu-n_\mu n_\nu)
 +Q(u_\mu n_\nu+n_\mu u_\nu).
\end{equation}
Here $T_{\hat0\hat1}=-Q$; $Q>0$ is outward energy flux. Then
\begin{equation}\label{eq:EFE}
 E_0=\kap\eps,\quad E_1=\kap P_r,\quad E_2=\kap P_\perp,\quad G_{tr}=-\kap ABQ.
\end{equation}
This decomposition permits flux and anisotropy; it does not supply their constitutive laws. Define
\begin{equation}
 \Dt=A^{-1}\partial_t,\quad\Dr=B^{-1}\partial_r,\quad U=\Dt R,\quad\Gamma=\Dr R,
 \quad H_r=\Dt\Lambda,\quad H_\perp=U/R,\quad a_r=\Dr\Phi.
\end{equation}
$\Theta=H_r+2H_\perp$ and $\sigma^2=(H_r-H_\perp)^2/3$. Projecting $\nabla_\mu T^{\mu\nu}=0$ along $u,n$ gives
\begin{align}
 \Dt\eps+(\eps+P_r)H_r+2(\eps+P_\perp)H_\perp+\Dr Q+2Q(\Gamma/R+a_r)&=0,\label{eq:energy}\\
 \Dt Q+2Q(H_r+H_\perp)+\Dr P_r+(\eps+P_r)a_r+2\Gamma(P_r-P_\perp)/R&=0.\label{eq:momentum}
\end{align}
Both are checked against an independent four-dimensional divergence. One must not use a no-flux Euler equation if retained spin correlations give $Q\neq0$.

\subsection{Constraints, mass and evolution identities}
Use geometric mass $M=Gm_{\rm MS}$, not baryonic mass:
\begin{equation}\label{eq:massdef}
 M=\frac R2(1+U^2-\Gamma^2).
\end{equation}
Differentiating this definition and substituting Eqs.~\eqref{eq:E0}--\eqref{eq:Etr} gives the exact mass identities
\begin{align}
 \Dr M&=\frac\kap2R^2(\eps\Gamma+QU), &
 \Dt M&=-\frac\kap2R^2(P_rU+Q\Gamma).\label{eq:masslaws}
\end{align}
The metric constraint and momentum constraint, with extrinsic-curvature convention $\mathsf K^i{}_j=-\operatorname{diag}(H_r,H_\perp,H_\perp)$, are
\begin{align}
 {}^{(3)}R+4H_rH_\perp+2H_\perp^2&=2\kap\eps,
 &{}^{(3)}R&=2(1-\Gamma^2)/R^2-4\Dr\Gamma/R,\\
 \Dr U-\Gamma H_r&=\kap RQ/2.\label{eq:ADM}
\end{align}
Here $\mathsf K$ denotes extrinsic curvature, not contortion $K^\rho{}_{\mu\nu}$. The radial and geometric evolution identities are
\begin{align}
 \Dt R&=U, &\Dt U&=\Gamma a_r-M/R^2-\kap RP_r/2,\label{eq:U}\\
 \Dt\Gamma&=Ua_r+\kap RQ/2.\label{eq:Gamma}
\end{align}
The angular equation \eqref{eq:E2} is an essential independent residual even when redundant on exact constrained matter evolution. None of Eqs.~\eqref{eq:masslaws}--\eqref{eq:Gamma} invents a shear or spin force: the source must be supplied by a supported matter model.

\subsection{FLRW and vacuum checks}
Set $\Phi=0$, $R=ar$, $B=a/\sqrt{1-kr^2}$. Then $U=HR$, $\Gamma=\sqrt{1-kr^2}$, $H_r=H_\perp=H$, and
\begin{equation}\label{eq:flrwEFE}
 3(H^2+k/a^2)=\kap\eps,\qquad -2\ddot a/a-H^2-k/a^2=\kap P,
 \qquad Q=0.
\end{equation}
Minkowski space $(A,B,R)=(1,1,r)$ and Schwarzschild exterior $(A^2,B^{-2},R)=(1-2M_0/r,1-2M_0/r,r)$ yield $G_{\mu\nu}=0$. The Schwarzschild component check uses a regular exterior chart; its coordinate singularity is not used as an interior fluid gauge.

\section{The three matter descriptions and their symmetry limits}
\subsection{Model 1: spinless perfect fluid}
Set the spin current and torsion to zero. A conserved-current action of the standard Brown type \cite{Brown} suffices to make the constitutive content explicit:
\begin{equation}\label{eq:fluidaction}
 I_{\rm pf}=\int\left[-\sqrt{-g}\,\eps(n,\chi)+J^\mu(\partial_\mu\varphi+\beta\partial_\mu\chi+\beta_A\partial_\mu\alpha^A)\right]\dd^4x,
 \quad n=\frac{\sqrt{-g_{\mu\nu}J^\mu J^\nu}}{\sqrt{-g}}.
\end{equation}
$J$ is a vector density; its potential enforces particle conservation, and $\beta$ enforces $u\cdot\nabla\chi=0$. Material labels allow general fluid variations; they introduce no additional stress when eliminated. At fixed current density $J$, $\delta n=\tfrac n2(g_{\mu\nu}+u_\mu u_\nu)\delta g^{\mu\nu}$ and $\delta\sqrt{-g}=-\tfrac12\sqrt{-g}g_{\mu\nu}\delta g^{\mu\nu}$. Substituting these variations into the material energy term gives the perfect-fluid metric source and $P=n\eps_{,n}-\eps$. Varying $J$ and the potentials gives
\begin{equation}
 \nabla_\mu(nu^\mu)=0,\quad \Dt\chi=0,\quad
 T_{\mu\nu}=(\eps+P)u_\mu u_\nu+Pg_{\mu\nu},\quad
 P=n\eps_{,n}-\eps.
\end{equation}
For $\eps=\rho=C(\chi)n^{1+w}$, $P=p=w\rho$ and $\tau=0$. Equations \eqref{eq:energy}--\eqref{eq:momentum} reduce to
\begin{equation}\label{eq:PFconservation}
 \Dt\rho=-(1+w)\rho\Theta,\qquad p'=-(\rho+p)\Phi'.
\end{equation}
For $\rho>0$, $0<w\leq1$ the fluid acoustic speed is $c_s^2=w$, real and causal; the full gravitational gauge and boundary problem must still be specified. Dust $w=0$ has coincident pressureless characteristics and requires a pre-caustic Lagrangian/dust formulation, not a strictly hyperbolic acoustic solver. For $w<0$ or $w>1$ the standard causal Euler interpretation fails; no such branch is authorized by this audit.

\subsection{Model 2: unpolarized, averaged Weyssenhoff fluid}
The microscopic source obeys a convective Frenkel condition, while
\begin{equation}
 \langle S^H_{\mu\nu}\rangle=0,\qquad
 S_H^2=\tfrac12\langle S^H_{\mu\nu}S_H^{\mu\nu}\rangle\ne0.
\end{equation}
Under local rotationally isotropic statistics in the rest frame,
\begin{align}
 \langle S^H_{ij}S^H_{kl}\rangle&=\frac{S_H^2}{3}(\delta_{ik}\delta_{jl}-\delta_{il}\delta_{jk}),\\
 \langle S^H_{i}{}^k S^H_{jk}\rangle&=\frac23 S_H^2\delta_{ij}.
\end{align}
We must first eliminate microscopic torsion and then average: $\langle K^2\rangle$ is not $\langle K\rangle^2$. Setting the mean torsion to zero at the start would erase the effect being modelled. The source used by the existing laboratory is
\begin{equation}\label{eq:avgsource}
 \eps=\rho-q,\quad P=w\rho-q,\quad q=\frac\kap4 S_H^2,\quad Q=0,\quad P_r=P_\perp=P.
\end{equation}
This follows the selected normalization in \cite{Hashemi}; the following \emph{additional} hypotheses are required for an inhomogeneous local use:
\begin{enumerate}
 \item Averaging is over a length small compared with resolved metric/matter variation, large compared with microscopic spin scales.
 \item Derivatives and averaging commute at retained order; correlations of fluctuating velocity, connection, acceleration and shear with spin are discarded or shown negligible. Vanishing mean spin alone is insufficient.
 \item The second moment is locally isotropic, and its statistics and spin per particle are transported without torque, relaxation, spin exchange or particle production.
 \item Metric backreaction/averaging commutators and unresolved anisotropic moments are absent at the approximation order.
\end{enumerate}
These hypotheses are constitutive and coarse-graining assumptions, not established by inserting a radial dependence in a homogeneous formula. Under them, a precise local effective action is Eq.~\eqref{eq:fluidaction} with
\begin{equation}\label{eq:effectiveEOS}
 \eps(n,\chi)=C(\chi)n^{1+w}-D(\chi)n^2,
 \qquad \rho=C(\chi)n^{1+w},\quad q=D(\chi)n^2.
\end{equation}
$D\geq0$ encodes conserved spin variance, not an arbitrary time-dependent repulsion. Variation yields $P=w\rho-q$ and
\begin{align}
 \Dt\rho&=-(1+w)\rho\Theta,&\Dt q&=-2q\Theta,&
 \Phi'&=-\frac{w\rho'-q'}{(1+w)\rho-2q}.\label{eq:avgtransport}
\end{align}
Equivalently $q/\rho^{2/(1+w)}$ is advected. Either $\rho,q$ or $n,\chi$ can be used, not both pairs as independent fluid degrees of freedom. Separate ordinary/spin conservation is part of the model; effective conservation alone would not determine two matter functions. The current action justifies the \emph{effective metric theory} under these hypotheses; it is not a substitute derivation of every microscopic Weyssenhoff action.

\subsubsection{Closure is not hyperbolicity}
Let $h=\eps+P=(1+w)\rho-2q$. At fixed advected composition, the acoustic principal symbol has
\begin{equation}\label{eq:sound}
 c_{\rm eff}^2=\left.\frac{\dd P}{\dd\eps}\right|_\chi
 =\frac{w(1+w)\rho-2q}{(1+w)\rho-2q}.
\end{equation}
Derive this by differentiating Eq.~\eqref{eq:effectiveEOS} with respect to $\ln n$. In a local inertial frame, the frozen-composition principal equations are
\begin{equation}
 \partial_t\delta\eps+h\partial_x\delta v=0,
 \qquad h\partial_t\delta v+\partial_x\delta P=0,
 \qquad \partial_t^2\delta\eps-c_{\rm eff}^2\partial_x^2\delta\eps=0.
\end{equation}
For $0\leq w<1$, the branch $h>0$ has $c_{\rm eff}^2<0$ whenever
\begin{equation}
 \frac{w(1+w)}2<\frac q\rho<\frac{1+w}2.
\end{equation}
In particular, ordinary dust with any $0<q<\rho/2$ has negative acoustic speed squared. Fourier modes grow as $\exp(|k|\sqrt{-c_{\rm eff}^2}t)$, without a wavelength-independent estimate: a generic smooth radial Euler initial-value problem is not well posed on this branch. At $h=0$ the usual Euler inversion degenerates. On the $h<0$ branch $c_{\rm eff}^2>1$ for $w<1$; this is a failure of a causal effective-fluid interpretation, not a theorem that the underlying EC action has propagating ghosts.

A flat homogeneous bounce has $q=\rho$, hence $c_{\rm eff}^2=w+2$. Its trajectory crosses the $h=0$ ratio $(1+w)/2$ before bouncing when these ratios are attained continuously. The spatially homogeneous ODE has no acoustic gradients and remains a consistent symmetry-restricted calculation; this observation does \emph{not} invalidate its existing numerical benchmark. It does prevent calling its unrestricted inhomogeneous continuation a verified causal PDE. For a positive-pressure low-spin regime, only states with $h>0$ and $0<c_{\rm eff}^2\leq1$ are candidates for a causal Euler implementation; approaching $c_{\rm eff}^2=0$ also requires a characteristic-degeneracy analysis. Artificial viscosity, a guessed pressure or changing spin transport is not an authorized cure.

\subsection{Model 3: selected full Weyssenhoff field formulation}
The phrase ``full Weyssenhoff'' does not select a unique averaging scheme or microscopic material action. We use the \emph{published field equations} \cite{BHL}, Eqs. (19)--(27), based on \cite{OK}. This is an allowed identification of an established formulation, not a newly manufactured action. Its full constrained matter-action derivation is an explicitly open reproducibility task.

Keep the BHL spin variable $b_{\mu\nu}$ for linear terms and use their Appendix B transformation to our metric signature: $g_B=-g$, $u_B^\mu=u^\mu$, $(u_B)_\mu=-u_\mu$, $(b_B)_{\mu\nu}=b_{\mu\nu}$. The resulting symmetric effective metric source and spin equation are
\begin{align}
 T_{\mu\nu}^{B}&=(\rho+p-2\kap b^2)u_\mu u_\nu+(p-\kap b^2)g_{\mu\nu}+\mathcal D_{\mu\nu},\label{eq:fullstress}\\
 \mathcal D_{\mu\nu}&=-2(g^{\alpha\beta}-u^\alpha u^\beta)\nabla_\alpha[u_{(\mu}b_{\nu)\beta}],\label{eq:Dspin}\\
 h_\mu{}^\alpha h_\nu{}^\beta\nabla_\lambda(u^\lambda b_{\alpha\beta})&=0,
 \qquad b_{\mu\nu}u^\nu=0,\quad h_\mu{}^\nu=\delta_\mu^\nu+u_\mu u^\nu.\label{eq:fulltransport}
\end{align}
The tensor in parentheses in Eq.~\eqref{eq:Dspin} is \emph{not} the spatial projector; replacing its minus sign by plus without transforming the source equation changes the model. We use $b^2=\tfrac12 b_{\mu\nu}b^{\mu\nu}$. The spin current in the BHL source convention is $u^\lambda b_{\mu\nu}$, their torsion is $Q_B=\kap u b$, and the operational source for Eq.~\eqref{eq:cartanform} is $-2u b$. This is why Eq.~\eqref{eq:dictionary} is needed. After elimination, $\nabla_\mu T_B^{\mu\nu}=0$ rather than a guessed conservation law for a bare canonical tensor is imposed.

Equation \eqref{eq:fulltransport} is equivalent to a projected spin-density transport equation. If $a_\mu=u^\alpha\nabla_\alpha u_\mu$, then
\begin{equation}
 u^\lambda\nabla_\lambda b_{\mu\nu}+\Theta b_{\mu\nu}
 =-2u_{[\mu}b_{\nu]\lambda}a^\lambda.
\end{equation}
Its contraction gives $\Dt b^2=-2\Theta b^2$ in this ideal nondissipative formulation. This is a published transport specialization; collision, relaxation or particle-production laws are not being supplied.

\subsubsection{Exact FLRW versus statistical FLRW}
A nonzero spatial two-form is dual to a nonzero spatial axial vector and selects a direction. Exact FLRW invariance of the spin/torsion configuration under the spatial isometry group forces $b_{\mu\nu}=0$ for this convective Frenkel source. An isotropic axial torsion three-form is a different possible EC source and must not be identified with $u^\lambda b_{\mu\nu}$. The coherent nonzero model therefore has \emph{no exact strong-symmetry FLRW realization}. This agrees with the symmetry warning in \cite{BB}. A statistically isotropic ensemble can have zero mean two-form but nonzero variance and an FLRW \emph{effective metric}; it then requires the averaging assumptions of Model 2 or a separately closed moment theory. There is no contradiction between these two statements of symmetry.

\paragraph{A verified metric-only exception.}
The preceding exclusion concerns the \emph{complete configuration}, not the metric alone. To prevent an overstatement, consider the selected ideal field formulation on a flat FLRW metric in Cartesian coordinates, with a spatially homogeneous aligned spin $b_{xy}=a^2\beta(t)$, all other spatial components zero. Its scalar is $b^2=\beta^2$, and direct substitution gives $\mathcal D_{\mu\nu}=0$ and $\dot\beta+3H\beta=0$. In spherical coordinates the same tensor has
\begin{equation}\label{eq:aligned}
 b_{r\phi}=a^2\beta r\sin^2\theta,\qquad
 b_{\theta\phi}=a^2\beta r^2\sin\theta\cos\theta.
\end{equation}
It is a smooth Cartesian aligned configuration, not a radial-spin profile with a singular center. The metric equations close as Eq.~\eqref{eq:homogeneous} with $q=\kap\beta^2$ and $\dot\beta=-3H\beta$. Thus a nonzero coherent spin field can be compatible with an \emph{FLRW metric} while failing rotational invariance of the full matter/torsion configuration. This is a symmetry-noninheriting source. The cancellation and its nonzero rotational Lie derivative are independently checked. It is an analytical compatibility example, not a proposed astrophysical spin preparation, and no classification of all such branches for $k\ne0$ is claimed. A generic ``full-spin FLRW'' mode would have to state which notion of symmetry and which formulation it uses.

\subsubsection{Exact spherical spin branch}
Under proper rotations about a regular center, a smooth covariant two-form has invariant basis $\dd t\wedge\dd r$ and $\sin\theta\,\dd\theta\wedge\dd\phi$. Frenkel with comoving $u$ kills the first. Thus the coherent ideal branch must have
\begin{equation}\label{eq:radspin}
 b_{\theta\phi}=R^2\sin\theta\,b(t,r),\qquad b_{\phi\theta}=-b_{\theta\phi},\qquad\text{all other components }0.
\end{equation}
This corresponds to a radial axial spin vector. It is allowed by $SO(3)$, but reflection/parity invariance of an axial radial field would force it to vanish. Smoothness at the center requires $b(t,r)=O(r)$ in a regular Cartesian extension, not a constant nonzero radial spin magnitude at $r=0$.

There is a particularly restrictive result. Substitute Eq.~\eqref{eq:radspin} into the retained derivative tensor \eqref{eq:Dspin}. Angular divergence vanishes; radial acceleration contracts with no spin component; angular shear eigenvalues coincide and their symmetric contraction with an antisymmetric angular spin tensor cancels. Hence
\begin{equation}\label{eq:radresult}
 \mathcal D_{\mu\nu}=0,\quad Q=0,\quad P_r=P_\perp=w\rho-\kap b^2,\quad\eps=\rho-\kap b^2,
 \quad\partial_t(BR^2b)=0.
\end{equation}
The cancellation is checked for arbitrary $\Phi,\Lambda,R,b$ by covariant differentiation, as is the Lie invariance under all three rotation generators. A nonspherical spin component produces nonzero derivative stress in a separate control test. The cancellation therefore does not silently discard terms: it follows from this precise symmetry-restricted ideal formulation. Tidal shear alone does not generate a thermodynamic $P_r\ne P_\perp$ here.

Equations \eqref{eq:radresult} plus $p=w\rho$ give a formally closed smooth branch on the same nonsingular-fluid domain as Model 2, with $q=\kap b^2$ and the same local acoustic obstruction. It remains distinct from an unpolarized ensemble. A regular nonzero coherent radial branch cannot be promoted to an exact full-FLRW spin configuration. If the proposed Model 3 instead retains nonzero averaged velocity--spin, shear--spin, spin-gradient or metric correlations, these require additional moment evolution and constitutive equations. Equations for arbitrary $Q,P_r,P_\perp$ in Section 2 are a source decomposition, not that missing closure. This broader proposal remains \status{not closed; implementation blocked}.

\section{Explicit systems, compatibility and independent data}
\subsection{The six proposed geometry--matter combinations}
\begin{longtable}{p{.13\textwidth}p{.19\textwidth}p{.32\textwidth}p{.25\textwidth}}
\toprule Geometry & Matter & Consistency/closure & Implementation decision\\\midrule\endhead
FLRW & 1: spinless $p=w\rho$ & Exact perfect-fluid limit; density constraint plus regular $(a,\dot a)$ evolution. & Supported homogeneous system; preserve existing module.\\
FLRW & 2: unpolarized variance & Effective metric symmetry only; requires local isotropic spin statistics and conserved spin variance. Algebraically closed ODE. & Supported as the documented Hashemi closure; no proof of radial stability or exterior matching.\\
FLRW & 3: coherent/full spin & Strong full-field FLRW symmetry excludes nonzero coherent spin. A verified flat metric-only aligned branch has non-isotropic spin/torsion; statistical variants need averaging closure. & Restricted declared-symmetry cases only; no generic full-spin switch.\\
Spherical & 1: spinless & Exact Einstein--Euler for $0<w\leq1$ in its smooth causal domain; dust pre-caustic case separately. Gauge and boundary data required. & Candidate first PDE stage after boundary/gauge acceptance tests; no PDE implemented here.\\
Spherical & 2: averaged variance & Conditional local effective action closes equations away from degeneracies. Dust low-spin branch has $c_{\rm eff}^2<0$; enthalpy can vanish. & General bounce-capable PDE not authorized. Restrict any future local test to a verified hyperbolic domain.\\
Spherical & 3: full formulation & Ideal coherent $SO(3)$ radial branch is formal and its linear-spin stress cancels. Parity, center and hyperbolicity restrictions apply. A broader retained-correlation model is not closed. & No generic full-Weyssenhoff PDE authorized; no invented anisotropic stress/transport.\\\bottomrule
\end{longtable}
This table concerns the FLRW and general spherical metrics requested here, not the separate retained Kantowski--Sachs explorer.

\subsection{Supported homogeneous equations}
For Model 1 set $q=0$; for Model 2 use $q=\kap S_H^2/4$. The independent regular system is
\begin{align}
 \dot a&=v,&\dot v&=-\frac\kap6 a[(1+3w)\rho-4q],\\
 \dot\rho&=-3(1+w)H\rho,&\dot q&=-6Hq,\qquad H=v/a,\label{eq:homogeneous}\\
 H^2+k/a^2&=\kap(\rho-q)/3.\label{eq:homconstraint}
\end{align}
The last equation is an initial constraint and monitored first integral, not an independent sign-selected square-root ODE. Matter and spin can be evaluated by their exact scaling laws rather than independently evolved twice. At fixed $a_i,k,w,q_i$, either $H_i$ or $\rho_i$ is constrained by the other. A positive-scale bounce satisfies $H=0$, $\dot H>0$ and $4q>\rho+3p$. Metric curvature and a continuation test are required in addition to an instantaneous event. For Model 3, strong full-field FLRW symmetry gives its zero-spin subcase. The flat metric-only aligned example in Eq.~\eqref{eq:aligned} instead uses $q=\kap\beta^2$, $\dot\beta=-3H\beta$ in this same metric system, while explicitly breaking isotropy of spin/torsion. A statistical reduction requires the assumptions of Model 2 or a new moment closure.

\subsection{Supported spherical perfect-source system}
For Model 1 $(\eps,P)=(\rho,w\rho)$ and $q=0$. For the \emph{conditional} Model 2 use Eq.~\eqref{eq:avgsource}; for the restricted coherent Model 3 use Eq.~\eqref{eq:radresult}. On a monotone shell chart $\Gamma\ne0$ with $h=\eps+P\ne0$:
\begin{align}
 \Dt R&=U,&\Dt U&=-\frac{\Gamma}{h}\Dr P-\frac{M}{R^2}-\frac\kap2 PR,\\
 \Dt M&=-\frac\kap2 PR^2U,&\Dt\Gamma&=-\frac U h\Dr P,\label{eq:closed}\\
 \Dt\rho&=-(1+w)\rho\left(\frac{\Dr U}{\Gamma}+\frac{2U}{R}\right),&
 \Dt q&=-2q\left(\frac{\Dr U}{\Gamma}+\frac{2U}{R}\right),\\
 M'&=\frac\kap2\eps R^2R',&\Gamma^2&=1+U^2-2M/R,\\
 B&=R'/\Gamma,&\Phi'&=-P'/h.
\end{align}
The displayed evolution equations include redundant quantities intentionally: one must choose a free-evolution or constrained scheme, not overwrite all evolved variables independently. In a constrained prototype, evolve $R,U,\rho,q$, solve $M'$ with a center condition, derive $\Gamma$ algebraically and $B$ kinematically, and solve the lapse-gradient constraint. Monitor the unused mass-time, $\Gamma$-time, angular Einstein and ADM identities. Alternatively evolve $M$ and monitor its radial constraint; the scheme specification must select one. The divided form above is forbidden at a neck $\Gamma=0$ or enthalpy zero. The undivided metric equations remain the starting point for any alternative formulation there.

For general retained flux/anisotropy, use Eqs.~\eqref{eq:EFE}, \eqref{eq:energy}--\eqref{eq:Gamma}, not Eq.~\eqref{eq:closed}. Unless a validated spin formulation supplies $Q,P_r,P_\perp$ and their transport, that is an underdetermined system. No source function is to be inferred merely from a desired bounce.

\subsection{Independent initial variables and gauge}
For a smooth monotone, comoving no-flux patch, independent profiles can be chosen as $R_i(r),U_i(r),\rho_i(r)$ plus $q_i(r)$ for Model 2. A radial relabelling can fix $R_i=r$; then Model 1 has two free profiles $(U_i,\rho_i)$ and Model 2 has three $(U_i,\rho_i,q_i)$, apart from constants, domain and boundary choices. For coherent Model 3 use signed $b_i$ rather than $q_i$ and impose center/parity conditions; $q_i=\kap b_i^2$ is not independent.

Integrate $M_i'=\kap\eps_iR_i^2R_i'/2$ with $M_i(0)=0$. Choose the outward branch
\begin{equation}
 \Gamma_i=+\sqrt{1+U_i^2-2M_i/R_i},\quad B_i=R_i'/\Gamma_i>0,
 \quad H_{r,i}=U_i'/R_i',\quad H_{\perp,i}=U_i/R_i.
\end{equation}
These solve the Hamiltonian/mass and momentum constraints; arbitrary independent choices of $B_i,M_i,H_{r,i}$ would over-specify initial data. In Model 2 require $\rho_i,q_i\geq0$, real geometry and a declared principal-symbol domain, not an unexamined positive effective density assumption. Integrate $\Phi_i'=-P_i'/h_i$ and fix a single lapse normalization, e.g. $\Phi(t,0)=0$. This sets central proper time, not all shell proper times.

The unknowns may be stored as $(R,U,M,\Gamma,B,\Phi,\rho,q)$, but $M,\Gamma,B,\Phi$ above are constraint/gauge variables in the constrained scheme. Number density/composition is an alternative representation. Spin current, scalar variance, effective density and geometric mass are different objects and must not be counted as additional independent data when derived. Canonical spin/torsion components are algebraic once the appropriate spin source is fixed.

Regular-center conditions include $R=\ell(t)r+O(r^3)$, $B=\ell(t)+O(r^2)$, $U=O(r)$, $M=O(r^3)$, $\Phi'=0$, and even smooth scalar density/variance profiles. Then $\Gamma=1+O(r^2)$ and $H_r-H_\perp=O(r^2)$. The coherent radial spin magnitude instead is odd to leading order, $b=O(r)$, so its variance vanishes at the center. Ensemble variance in Model 2 need not vanish there.

\subsection{Constraint propagation and clocks}
The contracted Bianchi identity, verified symbolically, enforces propagation of Einstein constraints provided the selected matter equations and consistent boundary/gauge conditions hold. This is not a stability proof of a discrete scheme, nor does it supply missing matter closure. Near $h=0$, division of a formally valid equation does not prove its continuation exists.

In general $\dd\tau_r=A(t,r)\dd t$: comoving shells do not share one proper-time coordinate when pressure gradients require $\Phi'\ne0$. Spherical plots must use a common declared coordinate time $t$ or explicitly record shell proper times. A future code must export $A(t,r)$, the lapse normalization and $\tau_r(t)$, rather than silently applying the FLRW common-proper-time assumption.

\section{Boundaries, exterior matching and observables}
\subsection{Finite matter region and spin-free exterior}
A source-free minimal EC exterior has zero torsion by the algebraic Cartan equation. With no other matter or cosmological constant it obeys the vacuum Einstein equation; a locally spherical vacuum region is Schwarzschild, subject to the usual domain assumptions. A Vaidya-type exterior instead contains a specified null radiative source and is not vacuum. Its mass function cannot be assigned from a desired interior plot.

For a timelike boundary $\Sigma$, a shell-free \emph{effective metric} matching requires continuity of the induced metric and Levi--Civita extrinsic curvature. For a comoving no-flux perfect source against vacuum, the normal traction conditions give
\begin{equation}\label{eq:matching}
 P|_\Sigma=0,\qquad M|_\Sigma=M_{\rm ext},\qquad\Dt M|_\Sigma=0.
\end{equation}
The last condition follows from the pressure-work equation when the first holds, not from identifying effective mass with conserved rest mass. Density may jump at a pressure-free dust boundary without necessarily creating a metric thin shell. The complete EC matching also involves the spin source and contortion. If there is a smooth spin-free collar, EC reduces to GR there and a conventional shell-free junction may be applied at an outer surface after all metric and matter conditions have been checked. Setting a plotted spin scalar to zero at one point does not construct such a collar or prove a valid junction.

For Model 2, $P=w\rho-q$. In ordinary dust, $P_\Sigma=-q_\Sigma$, so a nonzero homogeneous spin variance extending to a comoving vacuum interface cannot satisfy Eq.~\eqref{eq:matching}. A radial taper changes the matter problem and must be a constrained physical initial profile with compatible transport, not an imposed cosmetic spin window. At $q_\Sigma=0$, a $w>0$ fluid also needs $\rho_\Sigma=0$. Even in GR the constant-$w>0$ Euler lapse obeys
\begin{equation}
 \Phi=-\frac{w}{1+w}\ln\rho+f(t),
\end{equation}
so a finite comoving zero-density edge is not automatically a regular finite-lapse free boundary in this chart. A different surface equation of state, noncompact fluid region, specified atmosphere, moving interface or surface-layer model requires an explicit choice and analysis. No physical atmosphere or pressure floor is added here.

With a comoving boundary and no interior heat flux, matching to a genuinely radiating Vaidya exterior cannot be justified by a changing $M(t)$ alone. Continuity of normal energy flux and stress imposes additional conditions; an exterior radiation flux requires interior transport, a noncomoving interface with matter flux, or a surface layer. The sign and normalization of a pressure--flux condition must be derived for that chosen interface and null normalization, not copied across conventions.

For a full spinful interface, one must specify the induced coframe/metric, spin transport to the interface, jumps of connection/contortion, and whether a surface stress or surface spin current is allowed, then integrate the EC equations in an applicable junction formalism such as \cite{AKP}. A nonzero surface layer is not forbidden in principle, but its action/constitutive data and distributional regularity are absent here. In particular, a delta spin sheet cannot simply be squared in a torsion-squared source. The publisher abstract and bibliographic information for \cite{AKP} were accessible; a complete derivation of its component conditions was not independently reproduced. Thus \status{spinful surface-layer matching remains unresolved and blocks implementation of that branch}. No unsupported EC--Israel jump formula is supplied.

\subsection{Null expansions and mass}
With the outward normal aligned to increasing $r$ on the chosen star patch,
\begin{equation}
 \ell_\pm=u\pm n,\quad\ell_+\cdot\ell_-=-2,\qquad
 \theta_\pm=\frac2R(U\pm\Gamma),\qquad
 \theta_+\theta_- =\frac4{R^2}\left(\frac{2M}{R}-1\right).
\end{equation}
Both negative is future trapping; both positive is anti-trapping. A marginally outer trapped sphere has $\theta_+=0$, $\theta_-<0$. Compactness greater than one alone does not distinguish future from anti-trapping. A change of radial orientation also changes the coordinate labelling of outgoing and ingoing normals, so orientation must be recorded near a neck. These are local trapping diagnostics, not event horizons. A finite interior alone has no specified future null infinity or exterior causal completion from which an event horizon could be determined.

\subsection{Regular bounce, shell crossing and curvature}
For a noncentral Lagrangian shell, an areal bounce has $R>0$, $U=0$ and $\Dt U>0$; a recollapse reverses that sign. A radial proper stretch uses $B\,\dd r$, not $R'$ alone. Define the local comoving volume element and its scale by
\begin{equation}
 V_{\rm cell}\propto BR^2\dd r\dd\Omega,\quad
 a_V/a_{V,i}=(BR^2/B_iR_i^2)^{1/3},\quad \Dt\ln V_{\rm cell}=\Theta.
\end{equation}
A volume turning point has $\Theta=0$ and $\Dt\Theta>0$ for a minimum. It need not be an areal bounce. The vorticity-free Raychaudhuri equation is
\begin{equation}
 \Dt\Theta=-\Theta^2/3-2\sigma^2-\frac\kap2(\eps+P_r+2P_\perp)+\nabla_\mu a^\mu,
 \qquad \nabla_\mu a^\mu=\Dr a_r+2\Gamma a_r/R+a_r^2.
\end{equation}
For the averaged isotropic source, an instantaneous volume minimum therefore needs
\begin{equation}
 \kap(4q-\rho-3p)>4\sigma^2-2\nabla_\mu a^\mu\quad(\Theta=0).
\end{equation}
Unlike a geodesic homogeneous anisotropic toy model, an inhomogeneous pressured fluid also has the acceleration-divergence term. This local inequality alone proves neither regular continuation nor absence of shell crossings or a well-posed matter evolution.

Metric invariants can be computed from the warped geometry:
\begin{align}
 \Rsc&=R_h+2(1-X)/R^2-4\Box_hR/R,\\
 \mathcal K&=R_h^2+8H_{ab}H^{ab}/R^2+4(1-X)^2/R^4.
\end{align}
A Lorentzian contraction is not a positive-definite norm; do not assume every diagnostic invariant has a positive logarithm. Bounded samples of one scalar do not prove geodesic completeness. Check orthonormal curvature components in regular frames, fluid variables, invariants, convergence and continuation. Full EC model 3 additionally needs torsion/spin regularity; metric curvature is not the full connection curvature.

A noncentral $R\to0$ is a shell-focusing candidate. A zero of $R'$ with $R>0$ is not automatically shell crossing: a regular closed-FLRW equator in $R=a\sin\chi$, $B=a$ has $R'=0$ and finite curvature. In a no-flux mass equation, nonzero $M'$ with $R'\to0$ forces an effective-density divergence; $M'=0$ simultaneously requires a separate regular-neck/coordinate analysis. Track $B$, the comoving particle-volume Jacobian, neighbouring material shells and ordinary density, not just one denominator. A numerical chart stopping condition is not a proven physical singularity.

\subsection{Counts, e-folds and finite-interval statements}
Count minima and maxima separately for each material shell, exclude imposed initial turning data from collapse-generated bounce counts, and require a persistent sign reversal resolved under refinement. A complete oscillation requires an actually resolved minimum--maximum--minimum (or maximum--minimum--maximum) sequence with regular evolution and consistent boundary data throughout. Different shells can be out of phase; a single universal bounce count is not automatic.

Meaningful local e-folds are
\begin{equation}
 N_R=\ln(R/R_i),\qquad N_V=\frac13\ln(BR^2/B_iR_i^2),\qquad
 N_{\rm acc}=\int \frac\Theta3\,\mathbf1_{\Theta>0}\mathbf1_{\Dt^2a_V/a_V>0}\,\dd\tau_r,
\end{equation}
where $\Dt^2a_V/a_V=\Dt\Theta/3+\Theta^2/9$. $N_R$ measures angular areal expansion, not total volume; all three reduce appropriately to FLRW e-folds only in the homogeneous isotropic limit. Center values require regular limits rather than direct $0/0$ evaluation. These are kinematic measures, not matter entropy or inflation by definition. Neither a finite count of cycles nor finite-interval expansion establishes eternal oscillation or eternal expansion.

\section{Validation, literature audit and implementation decision}
\subsection{Independent checks performed}
The reproducible Python suite builds the Levi--Civita connection and four-dimensional Ricci tensor directly from the diagonal spherical metric. It compares them to the Hessian/warped-product derivation and explicit components above. Further checks cover spherical stress conservation including flux and anisotropic pressures, geometric mass laws including flux, ADM constraints and evolution identities, contracted Bianchi identities, and the exact FLRW reduction for arbitrary $k$. Minkowski and Schwarzschild have zero Einstein tensor. Direct coordinate Riemann contraction gives $\mathcal K=48M_0^2/R^6$ in a Schwarzschild exterior test.

An inhomogeneous parabolic LTB dust solution
\begin{equation}
 R(t,r)=\left(\frac92M(r)\right)^{1/3}[t_s(r)-t]^{2/3},\quad B=R',\quad\Phi=0,
 \quad\kap\rho=\frac{2M'}{R^2R'}
\end{equation}
checks nontrivial spherical GR without pressure; test profiles are manufactured analytical validation data, not EC spin profiles. The tests also read the \emph{unchanged} existing FLRW simulator through a Node subprocess and compare its normalized matter/acceleration samples to the spherical Einstein equations. Separate symbolic checks verify contortion, the auxiliary torsion coefficient, isotropic covariance contraction, the effective EOS derivative, spin transport, coherent radial rotation invariance, and cancellation of radial linear-spin stress with a nonspherical nonzero control case. A flat FLRW metric with homogeneous aligned spin is checked as a distinct symmetry-noninheriting example. Acoustic sign and boundary-pressure obstructions are analytical checks rather than a time-dependent PDE simulation.

These tests validate algebra and stated special cases. They do not validate microscopic averaging, existence through an enthalpy-zero surface, general hyperbolicity of a full spin-moment system, a spinful junction construction or a global black-hole spacetime. The test inventory and exact reproduction commands are in \texttt{README.md}; the implementation gates are in \texttt{PDE\_SPEC.md}. No symbolic expression dump is used in place of the geometric derivation.

\subsection{Literature comparison and apparent contradictions}
\begin{longtable}{p{.18\textwidth}p{.34\textwidth}p{.37\textwidth}}
\toprule Reference & What is taken from it & What cannot be transferred automatically\\\midrule\endhead
Trautman \cite{Trautman} & Minimal EC variational structure, algebraic torsion and importance of source definitions. & A generic matter Lagrangian does not define a unique Weyssenhoff EOS, averaging rule or junction model.\\
BHL \cite{BHL}, based on \cite{OK} & Specific full effective stress and spin transport; signature and connection-slot dictionary. & Their $\kap b^2$ and linear-spin stress cannot be replaced numerically by Hashemi's named spin variable without normalization/sign conversion. Their specialized constraint discussion is not a proof of the proposed spherical PDE's well-posedness.\\
Hashemi et al. \cite{Hashemi} & The laboratory's averaged coefficient, conserved-spin homogeneous source and collapse benchmark. & Homogeneous ODE regularity has no radial acoustic principal symbol. A finite homogeneous pressure boundary is not automatically a valid vacuum junction.\\
B\"ohmer--Bronowski \cite{BB} & Strong cosmological symmetry acts on torsion/spin as well as the metric. & Excluding a coherent exact FLRW spin tensor does not exclude a statistically isotropic effective metric sourced by spin variance.\\
Luz--Mena--Ziaie \cite{LMZ} & An inhomogeneous EC effective-fluid derivation and explicit emphasis on source assumptions. & Their effective dust imposes $p-q=0$, with a quadratic ordinary-pressure relation in the discussed material model. It is \emph{not} the ordinary $p=0$, $P=-q$ conserved-spin closure tested here. Their non-singularity restrictions do not remove Eq.~\eqref{eq:sound}'s obstruction for that different closure.\\
Pop\l awski \cite{Pop} & Spherical Tolman geometry and a relativistic spin-fluid/thermodynamic setting. & A synchronous comoving metric requires absent pressure gradients (or an equivalent geodesic restriction). Particle production changes the number/spin evolution and is excluded here; its bounce/cycle claims cannot be imported into fixed-$w$, conserved-spin generic inhomogeneity.\\
AKP \cite{AKP} & Identified foundational EC junction reference. & The full spinful distributional matching calculation has not been reproduced; a publisher abstract is not a checked set of component jump conditions.\\\bottomrule
\end{longtable}
No novelty is claimed. The sound-speed calculation and spherical radial-spin reduction are audit deductions under the displayed assumptions. They should be reviewed against a selected microscopic fluid action and independent mathematical analysis before publication as a general physical result.

\subsection{Decisions and unresolved items}
\begin{enumerate}
 \item \status{Verified geometry and GR reduction:} a spinless spherical prototype may be specified, with gauge, boundary problem, characteristic treatment and center regularization established first. This task supplies a specification only.
 \item \status{Conditional averaged EC equations:} closed effective equations are written, but no generic dust or bounce-through-enthalpy-zero PDE is authorized. The acoustic obstruction is a demonstrable mathematical reason to withhold that implementation, not a reason to change the existing FLRW ODE.
 \item \status{Specific full ideal branch:} a coherent radial $SO(3)$ branch and its restrictions are identified. It does not supply the desired general correlation-retaining source and retains the same effective-fluid obstruction when specialized to the ideal EOS.
 \item \status{Open matter physics:} microscopic EOS at high density, collision/relaxation and moment closure, averaging/derivative commutators, coherent polarization, interface spin transport and an independently reconstructed full matter action remain separate obligations. No new laws have been invented to fill them.
 \item \status{Open boundary physics:} finite linear-EOS vacuum interfaces, full spinful surface layers, radiative matching and global exterior causal completion must be chosen and derived. A spin-free vacuum exterior is not by itself a boundary condition for the interior.
 \item \status{Open numerical analysis:} discrete constraint propagation, a well-posed gauge and characteristic system, shocks and material caustics, hyperbolicity-domain exit, validated error control and matching residual tolerances must precede production PDE outcomes.
\end{enumerate}
The appropriate next step is review of these gates, not a new ``general spherical EC bounce'' button. Existing runtime code, numerical solvers, FLRW exports and the GitHub Pages entry point remain unchanged.

\begin{thebibliography}{20}
\bibitem{Trautman} A. Trautman, Einstein--Cartan theory (2006), \href{https://arxiv.org/abs/gr-qc/0606062}{gr-qc/0606062}, especially the variational formulation and spin source conventions.
\bibitem{BHL} S. D. Br\'echet, M. P. Hobson and A. N. Lasenby, Weyssenhoff fluid dynamics in general relativity using a 1+3 covariant approach, Class. Quantum Grav. 24 (2007), \href{https://arxiv.org/abs/0706.2367v2}{arXiv:0706.2367v2}, Eqs. (1)--(5), (19)--(28), (60)--(63), Appendix B.
\bibitem{Hashemi} M. Hashemi, S. Jalalzadeh and A. H. Ziaie, Collapse and dispersal of a homogeneous spin fluid in Einstein--Cartan theory (2015), \href{https://arxiv.org/html/1407.4103v2}{arXiv:1407.4103v2}, Eqs. (2), (5), (8)--(9), (16)--(22).
\bibitem{OK} Yu. N. Obukhov and V. A. Korotky, The Weyssenhoff fluid in Einstein--Cartan theory, Class. Quantum Grav. 4 (1987) 1633--1657, \href{https://doi.org/10.1088/0264-9381/4/6/021}{doi:10.1088/0264-9381/4/6/021}. The published formulation as reproduced in \cite{BHL} is used here; its complete constrained matter-action derivation is not independently reconstructed.
\bibitem{BB} C. G. B\"ohmer and P. Bronowski, The homogeneous and isotropic Weyssenhoff fluid, \href{https://arxiv.org/abs/gr-qc/0601089v2}{gr-qc/0601089v2}; distinction between metric symmetry and symmetry of the complete torsion configuration.
\bibitem{Brown} J. D. Brown, Action functionals for relativistic perfect fluids, Class. Quantum Grav. 10 (1993) 1579, \href{https://arxiv.org/abs/gr-qc/9304026}{gr-qc/9304026}; conserved-current material action.
\bibitem{MS} C. W. Misner and D. H. Sharp, Relativistic equations for adiabatic, spherically symmetric gravitational collapse, Phys. Rev. 136 (1964) B571, \href{https://doi.org/10.1103/PhysRev.136.B571}{doi:10.1103/PhysRev.136.B571}.
\bibitem{LMZ} P. Luz, F. C. Mena and A. H. Ziaie, Influence of intrinsic spin in the formation of singularities for inhomogeneous effective dust space-times, Class. Quantum Grav. 36 (2019) 015003, \href{https://arxiv.org/abs/1811.12292v1}{arXiv:1811.12292v1}, especially Eqs. (17)--(24), (31). Their effective dust is not ordinary dust with a negative spin pressure.
\bibitem{Pop} N. Pop\l awski, Gravitational collapse of a fluid with torsion into a universe in a black hole, JETP 132 (2021) 374, \href{https://arxiv.org/html/2008.02136v2}{arXiv:2008.02136v2}. The version inspected is the 20 March 2026 revision; particle production and thermodynamic assumptions are not imported here.
\bibitem{AKP} W. Arkuszewski, W. Kopczy\'nski and V. N. Ponomariev, Matching conditions in the Einstein--Cartan theory of gravitation, Commun. Math. Phys. 45 (1975) 183--190, \href{https://doi.org/10.1007/BF01629248}{doi:10.1007/BF01629248}. Bibliographic/publisher abstract verified; a complete component-level reproduction of its junction formalism is not claimed.
\end{thebibliography}
\end{document}
